% ---- pg-159 （书页 147） ----

\begin{mvfigure}
\mvinc{pg-159}{0.189,0.124,0.849,0.550}
\phantomsection\label{fig:4.11}\mvcap{图 4.11\quad 木板硬度数据的散布图,其中样本 9 和 16 用实圆点代表}
\end{mvfigure}
\hfill$\blacksquare$

如果离群值被确定, 应当对其内容细查, 正如例 4.15 中在木板硬度数据的例子中所做的. 视离群值的特性和研究的目的而定, 在以后的分析中, 离群值可以删去或适当地“加权”.

虽然许多统计方法假定正态总体, 那些基于样本均值向量的方法通常不会受几个中等的离群值干扰. 霍金斯(Hawkins)(参见文献[7])给出了关于离群值问题的广泛论述.

\section{变换到接近正态性}  % 原书 4.8
\phantomsection\label{sec:4.8}

如果正态性不是一个可行的假定, 下一步该怎么办? 一个选择是不管正态性检验的结论, 把数据当做正态分布来继续进行. 我们不推荐这种做法. 因为在许多情况下, 它可能会导致错误的结论. 第二个选择是通过考虑数据变换, 使非正态数据更加“朝向正态”. 在适当数据变换后, 就可以实现正态理论分析.

变换无非是用不同的单位对数据进行重新表述. 例如, 正观测值的直方图显示出长的右尾时, 通过取对数或平方根来变换观测值, 经常显著地改进其关于均值的对称性, 并且提高对正态分布的近似程度. 新单位对所研究的特性常常提供更自然的表达式.

按照(1)理论上的考虑或(2)数据本身(或两者)提出适当的变换. 理论上已经证明对属于计数的数据取平方根通常可使之更接近正态. 类似地, 对数变换适用于比例, 而将费希尔(Fisher)的 $z$ 变换应用于相关系数, 可产生近似正态分布的量.
